This study compares AdaBoost, XGBoost, and LightGBM for binary healthcare risk prediction on three datasets with different sizes, feature representations, and class prevalences. A common protocol is used: stratified 80/20 development--test separation, persisted stratified five-fold partitions, fold-local preprocessing, fixed model configurations, and an untouched hold-out evaluation. Predictive performance is measured with ROC-AUC, PR-AUC, Recall, and F1-score, with PR-AUC as the primary ranking metric.

Under this protocol, XGBoost achieved the best overall cross-dataset rank (mean rank $1.0000$), with mean ROC-AUC $0.8844$, PR-AUC $0.5730$, and Recall $0.8202$. AdaBoost achieved the highest cross-dataset mean F1-score ($0.5392$), showing that the preferred model depends on the evaluation objective. Predictive variability was substantially higher on the smaller Heart Failure Prediction dataset: for example, XGBoost's ROC-AUC standard deviation was $0.0332$, compared with $0.0031$ and $0.0026$ on the two larger datasets.

For explanation stability, fold-level SHAP rankings were compared using Kendall's $\tau$, Spearman's $\rho$, and top-10 Jaccard similarity. Across the nine dataset--model combinations, mean Spearman correlation ranged from $0.9324$ to $0.9921$, while top-10 Jaccard similarity ranged from $0.7758$ to $1.0000$. These results indicate generally consistent explanations, although complete-ranking agreement does not always imply identical membership of the ten most important features.

Overall, the study shows that XGBoost was strongest under the selected ranking-oriented protocol, but no model dominated every metric and dataset. The contribution is a controlled comparison that evaluates predictive strength, predictive stability, and explanation stability together; the conclusions remain limited to the three datasets, fixed configurations, and evaluation protocol used here.
